

\Bsec{Interpolation}{
		
		\begin{itemize}
			\itemf Basisfunktionen
			\itemf Newton-Verfahren vs. Aitken-Neville
			\itemf Hermite und Spline-Interpolation
			\itemf Verständnis/Trick Fragen
		\end{itemize}
		}
		
%----------------------------------------------------------------------------------
\bsubsec{Basisfunktionen-Polynominterpolation}{
%		
%		Polynom-Interpolation mit Basisfunktionen $g_0$ bis $g_n$:
%		$$\begin{bmatrix}
%		g_0(x_0) & g_1(x_0) & \dots & g_n(x_0)\\
%		g_0(x_1) & \dots & \dots & g_n(x_1)\\
%		g_0(x_n) & \dots & \dots & \dots
%		c_0 \\ \dots \\ c_{n-1} \\ c_n
%		y_0 \\ \dots \\ y_{n-1} \\ y_n
%		$$ 
%		Lagrange-Basis:
%		$$L_j = \prod_{i = 0; i \neq j}^n \frac{x - x_i}{x_j - x_i}$$
%		
%	}
%		
		Stelle mit folgenden Stützstellen und Basisfunktionen ein Interpolationspolynom auf:
		
		\begin{enumerate}
			\itemf Mit den ''normalen'' Polynom-Basisfunktionen:
				$$f(1) = 1 \aae f(2) = 3 \aae f(3) = 5$$
			\itemf Mit den Lagrange-Basisfunktionen:
				$$P_0(0|0) \aae P_1(2|5) \aae P_2(4|6)$$
			\itemf $t(-2) = -3 \aae t(-1) = 1 \aae t(0) = 1 \aae t(1) = 3$ mit folgenden Basisfunktionen:
				$$g_0(x) = 1 \aae g_1(x) = x \aae g_2(x) = 2x^2 - 2 \aae g_3(x) = 4x^3 - 3x$$
			\itemf $x_0 = 0 \aae x_1 = 4 \aae y_0 = 2 \aae y_1 = 4$ mit den folgenden Basisfunktionen
			$$g_0(x) = x + 1 \aae g_1(x) = x^2$$
		\end{enumerate}
		
	}
	\bsol{
		
		\begin{enumerate}
			\itemf $p_0(t) = 2x - 1$
			\itemf $p_1(t) = \frac{7}{2}x - \frac{1}{2}x^2$
			\itemf $p_2(t) = x^3 + x^2 + 1$
			\itemf Ist Quatsch, bei Basisfkten mit Stützstellenanzahl $n$ kann man nur $(n - 1)$ als maximalen Polynomgrad erhalten, die $x^2$ ist 2. Grades und wir haben nur zwei Stützstellen! \n 
			Da würde das Polynom $2x + 2 - \frac{3}{8}x^2$ herauskommen, obwohl es nach korrekter Polynominterpolation $\frac{1}{2}x + 2$ ist.
		\end{enumerate}
	}
%----------------------------------------------------------------------------------
\bsubsec{Newton-Verfahren vs. Aitken-Neville}{
%		
%		Newton: $$f_{(i, k)} = \frac{f_{i+1, k-1} - f_{i, k-1}}{x_{i+k} - x_{i}}$$
%				$$[neu]f = \frac{f_{linksUnten} - f_{links}}{x_{ganzLinksUnten} - x_{links}}$$
%		Aitken-Neville: $$p[i,k] := p[i, k-1] + \frac{x-x[i]}{x[i+k]-x[i]} \cdot (p[i+1, k-1] - p[i, k-1])$$
%				$$p_{neu} := p_{links} + \frac{x-x_{links}}{x_{ganzLinksUnten}-x_{links}} \cdot (p_{linksUnten} - p_{links})$$
%	}
%		
%		Berechne das Interpolationspolynom mithilfe des Newton-Verfahrens und werte die zweite Aufgabe direkt an der gegebenen Stelle mithilfe von Aitken-Neville aus. 
		\begin{enumerate}
			\itemf Berechne das Interpolationspolynom mithilfe des Newton Verfahrens und folgenden Stützstellen: $$P_0(0|3) \aae P_1(1|2) \aae P_2(2|1) \aae P3(3|1)$$
			\itemf Werte mit dem Aitken-Neville Verfahren direkt an der Stelle $x = -1$ aus, mit den Stützstellen $$P_0(0|0) \aae P_1(1|2) \aae P(2|12) \aae P(-3|-18)$$
		\end{enumerate}
		
	}
	\bsol{
		
		\begin{enumerate}
			\itemf $$p(x) = \frac{1}{6}x^3 - \frac{1}{2}x^2 - \frac{2}{3}x + 3$$
			\itemf $$p(-1) = 0$$
		\end{enumerate}
	}

%----------------------------------------------------------------------------------
\bsubsec{Hermite und Spline-Interpolation}{
%	
%	$$p(t) = a_3 \cdot t^3 + a_2 \cdot t^2 + a_1 \cdot t + a_0$$
%	Allgemeine Hermite-Funktion:
%	$$p_i(t_i(x)) = y_i\cdot H_0(t_i(x))+y_{i+1}\cdot H_1(t_i(x))+h_i\cdot y_i'\cdot H_2(t_i(x))+h_i\cdot y_{i+1}'\cdot H_3(t_i(x))$$
%	Allgemeine Transformationsfunktion:
%	$$t_i(x) := \frac{x-x_i}{x_{i+1} - x_i} \in [0;1]$$
%	LGS von Splines:
%	$$\begin{bmatrix}
%	4 & 1 \\
%	1 & 4 & ...\\
%	& ... & ... & 1\\
%	&   & 1 & 4
%	y_1' \\ y_2' \\ ... \\ y_{n-2}' \\ y_{n-1}'
%	y_2 - y_0 \\
%	y_3 - y_1 \\
%	... \\
%	y_{n-1} - y_{n-3} \\
%	y_n - y_{n-2}
%	y_0' \\ 0 \\ ... \\ 0 \\ y_n'
%}
%	
	Gegeben sei: \\
	\begin{table}[H]
		\label{table:splinetask}
		\centering
		\begin{tabular}{cccccc}
			\toprule
			$i$ 	& 0 	& 1 	& 2 	& 3		& 4 \\ \midrule
			$x_i$	& -2  	& 0 	& 2 	& 4		& 6 \\
			$y_i$	& -3 	& -1 	& 1    	& 3  	& $\nicefrac{13}{3}$  \\
			$y'_i$  & -14  	& ?   	& ?    	& ?     & 1\\
			\bottomrule
		\end{tabular}
	\end{table}
	Außerdem sind die kubischen Basispolynome $H_0(t) \dots H_3(t)$ schon korrekt im Intervall $[0,1]$ ausgerechnet worden.

	\begin{enumerate}
		\itemf Berechne die fehlenden Ableitungen!
	\end{enumerate}
	
}
\bsol{
	
	\begin{enumerate}
		\item 
		$$\begin{bmatrix}
		4 & 1 & 0\\
		1 & 4 & 1\\
		0 & 1 & 4\\
		\end{bmatrix} \cdot
		\begin{bmatrix}
		y_1' \\ y_2' \\ y_3'
		\end{bmatrix} = \frac{3}{2} \cdot
		\begin{bmatrix}
		4 \\
		4 \\
		\nicefrac{10}{3}
		\end{bmatrix} - 
		\begin{bmatrix}
		-14 \\ 0 \\ 1
		\end{bmatrix}$$
		$$\begin{bmatrix}
			4 & 1 & 0 &|& 20\\
			1 & 4 & 1 &|& 6\\
			0 & 1 & 4 &|& 4
		\end{bmatrix} \longrightarrow
		\begin{bmatrix}
			4 & 1 & 0 &|& 20\\
			0 & 15 & 4 &|& 4\\
			0 & 0 & 56 &|& 56
		\end{bmatrix}$$
		$$y_1' = 5 \aae y_2' = 0 \aae y_3' = 1$$
	\end{enumerate}
}

\btask{
	
	\begin{table}[H]
		\label{table:splinetask2}
		\centering
		\begin{tabular}{cccccc}
			\toprule
			$i$ 	& 0 	& 1 	& 2 	& 3		& 4 \\ \midrule
			$x_i$	& -2  	& 0 	& 2 	& 4		& 6 \\
			$y_i$	& -3 	& -1 	& 1    	& 3  	& $\nicefrac{13}{3}$  \\
			$y'_i$  & -14  	& 5   	& 0    	& 1     & 1\\
			\bottomrule
		\end{tabular}
	\end{table}
	\begin{enumerate}
		\item Wieviele Splines definieren wir, in welchem Intervall und mit welchen Transformationsfunktionen?
	\end{enumerate}
	
}

\bsol{
	\begin{enumerate}
		\item Wir erhalten vier einzelne Splines, alle mit einer Intervallbreite von zwei. Die Transformationsfunktionen (Index jeweils identisch zur dazugehörigem Spline):
		$$t_0(x) = \frac{x+2}{2} \aae x\in [-2;0] \aae t_0(x) \in [0;1]$$
		$$t_1(x) = \frac{x}{2} \aae x\in [0;2] \aae t_1(x) \in [0;1]$$
		$$t_2(x) = \frac{x-2}{2} \aae x\in [2;4] \aae t_2(x) \in [0;1]$$
		$$t_3(x) = \frac{x-4}{2} \aae x\in [4;6] \aae t_3(x) \in [0;1]$$
	\end{enumerate}
}

\btask{
	
	\begin{table}[H]
		\label{table:splinetask3}
		\centering
		\begin{tabular}{cccccc}
			\toprule
			$i$ 	& 0 	& 1 	& 2 	& 3		& 4 \\ \midrule
			$x_i$	& -2  	& 0 	& 2 	& 4		& 6 \\
			$y_i$	& -3 	& -1 	& 1    	& 3  	& $\nicefrac{13}{3}$  \\
			$y'_i$  & -14  	& 5   	& 0    	& 1     & 1\\
			\bottomrule
		\end{tabular}
	\end{table}
	
	\begin{enumerate}
		\item Stelle nun die fertige Spline-Funktion $s(x)$ auf
		\item Definiere die einzelnen Splines mithilfe der kubischen Basispolynome.
	\end{enumerate}
	
}
\bsol{
	
	\begin{enumerate}
		\item $$s(x) = \begin{Bmatrix}
		p_0(t_0(x)) \aae x\in[-2;0] \aae t_0(x) = \frac{x+2}{2}\\
		p_1(t_1(x)) \aae x\in[0;2]  \aae t_1(x) = \frac{x}{2}\\
		p_2(t_2(x)) \aae x\in[2;4]  \aae t_2(x) = \frac{x-2}{2}\\
		p_3(t_3(x)) \aae x\in[4;6]  \aae t_2(x) = \frac{x-4}{2}
		\end{Bmatrix}$$
		\item 
		$$p_0(t_0(x)) = (-3)\cdot H_0(t_0(x))+ (-1)\cdot H_1(t_0(x))+ (-28)\cdot H_2(t_0(x))+ 10\cdot H_3(t_0(x))$$
		$$p_1(t_1(x)) = (-1)\cdot H_0(t_1(x))+ 1\cdot H_1(t_1(x))+ 10\cdot H_2(t_1(x))+ 0\cdot H_3(t_1(x))$$
		$$p_2(t_2(x)) = 1\cdot H_0(t_2(x))+ 3\cdot H_1(t_2(x))+ 0\cdot H_2(t_2(x))+ 2\cdot H_3(t_2(x))$$
		$$p_3(t_3(x)) = 3\cdot H_0(t_3(x))+ \frac{13}{3}\cdot H_1(t_3(x))+ 2\cdot H_2(t_3(x))+ 2\cdot H_3(t_3(x))$$
	\end{enumerate}
}

%	
%	
%	
%}
%	
%}

%----------------------------------------------------------------------------------
\bquestion{
	}
	\btask{
		
		Beantworte folgende (Trick-)Fragen möglichst genau:
		\begin{enumerate}
			\item Wie unterscheidet sich (ganz allgemein) das Ergebnis einer Interpolation mit den Lagrange-Basisfunktionen mit dem eines Newton-Verfahrens?
			\item Mehr Stützstellen, höhere Genauigkeit des Interpolationspolynoms?
			\item Wir möchten eine Interpolation durchführen, bei der wir relativ einfach neue Stützstellen hinzufügen können, welches Verfahren nutzen wir?
			\item Welchen speziellen Vorteil (außer Genauigkeit) bietet stückweise Interpolation gegenüber ''normaler''?
			\item Wir möchten das Intervall $[-3;7]$ auf das Intervall $[0;1]$ transformieren. Wie lautet die Transformationsfunktion $t(x)$?
		\end{enumerate}
		
	}
	\bsol{
		
		\begin{enumerate}
			\item Gar nicht, Interpolationspolynome sind eindeutig und beide können bei gleicher Stützstellenanzahl denselben Grad eines Polynoms erreichen.
			\item Theoretisch ja, aber bei zuvielen Stützstellen kann der ''Runge-Effekt'' auftreten und das Ergebnis wird wieder deutlich ungenauer.
			\item Wir nutzen das Newton-Verfahren, da man der Tabelle unten problemlos neue Stützstellen hinzufügen und netterweise noch relativ schnell ist.
			\item Es wirkt sehr stark dem ''Runge-Effekt'' entgegen.
			\item $$t(x) = \frac{x + 3}{10}$$
		\end{enumerate}
	}
%----------------------------------------------------------------------------------